\section{Executive synthesis}

\subsection{The problem is an interface problem}

A proof in intuitionistic primitive recursive arithmetic is not obtained merely by observing that a familiar argument is finite ``in principle.'' Four distinct claims must be separated.

\begin{enumerate}
\item \textbf{Representability.} Every object manipulated by the argument must have a fixed natural-number coding, and every operation used by the argument must belong to a fixed primitive-recursive family.
\item \textbf{Quantitative adequacy.} Whenever the ordinary theorem invokes convergence, compactness, continuity, a finite cover, an isolated branch, or a height bound, the formal statement must expose the numerical modulus or certificate actually consumed by the construction.
\item \textbf{Constructive coherence.} Finite witnesses obtained at different precisions need not agree. One must either prove that all valid choices have the same observable effect or build a compatible finite path by bounded search. A hidden appeal to an infinite branch, an exact sign, or a canonical label is not acceptable.
\item \textbf{Object-language reflection.} A correct Lean function, Python program, or external finite proof does not by itself yield an iPRA theorem. The end product must be an actual derivation of type
\[
   \Gamma \;\vdash[\PRAAx]\; \varphi,
\]
that is, a value of the repository's inductive relation \(\Prf\,\PRAAx\).
\end{enumerate}

The repository already addresses the first three points in many individual packets and several general frameworks. The most persistent cross-packet gap is the fourth. The central contribution of the present package is therefore a proof-producing interface whose semantic and syntactic sides meet at a small, explicit theorem: finite-trace reflection by quantifier-free induction.

\subsection{The five-layer architecture}

Every formalization target is divided into five layers. The division is not cosmetic: it is the main safeguard against accidentally treating a host-language theorem as an object-theory proof.

\paragraph{Layer I: fixed presentation.}
A presentation is a meta-level package of fixed primitive-recursive symbols. It may describe finite data exactly or infinite data by finite observations. It is deliberately not a natural-number code for an arbitrary program. This avoids the impossible demand for a total primitive-recursive universal evaluator.

\paragraph{Layer II: quantitative contract.}
A construction advertises the precision, search, guard, and certificate data it needs. Its contract consists of a primitive-recursive demand map, a bounded verifier, and an observable postcondition. A producer may be included, but a verifier is meaningful even when the difficult producer remains an open mathematical obligation.

\paragraph{Layer III: finite execution.}
The construction is executed as a bounded trace, a finite contract DAG, a bounded search, or a finite relational dynamic program. Every loop has an explicit primitive-recursive bound; every branch is Boolean or remains relational; every finite quantifier has a displayed bound.

\paragraph{Layer IV: object proof.}
The local base, step, and final contracts are themselves \(\Prf\,\PRAAx\) derivations. A generic compiler combines them using substitution, cut, the defining equations of primitive recursion, and \code{PRA.qfInd}. This layer contains no appeal to the intended standard model.

\paragraph{Layer V: transport.}
Only after the ordinary iPRA derivation is complete is it transported to the small native scalar core by the already checked theorem \code{pra_iff_native}. Keeping transport last makes the trust boundary visible.

\subsection{A single normal form for many proofs}

A large class of packet arguments can be reorganized into the following statement. Fix input \(x\), requested precision \(q\), and a bounded state evolution
\[
  s_0=I(x,q),\qquad s_{i+1}=S(i,s_i,x,q),\qquad i<N(x,q).
\]
Let \(\Phi(i,s,x,q)\) be a quantifier-free invariant after bounded-formula compilation. Suppose iPRA proves the three local obligations
\begin{align*}
  \Good(x)&\to \Phi(0,I(x,q),x,q),\\
  \Good(x)\wedge \Phi(i,s,x,q)&\to \Phi(i+1,S(i,s,x,q),x,q),\\
  \Good(x)\wedge \Phi(N(x,q),s,x,q)&\to \Accept(s,x,q).
\end{align*}
Then the primitive-recursive run symbol \(\Run\) satisfies
\[
  \iPRA\vdash \Good(x)\to
  \Accept(\Run(N(x,q),x,q),x,q).
\]
The proof is a direct application of quantifier-free induction to
\(
  \Phi(i,\Run(i,x,q),x,q).
\)
The recursion equations rewrite the zero and successor cases. This is elementary mathematics, but it is exactly the elementary mathematics that turns a bounded algorithm into an object-language proof.

The same normal form handles matrix elimination, finite search, interval bisection, finite-net minimization, construction of a simplicial chain map, bounded descent, certificate checking, and the forward pass of a contract DAG. A complicated theorem may still require difficult estimates, but once those estimates are presented as finite contracts, the proof-engineering problem becomes uniform.

\subsection{What is genuinely resolved}

The framework gives complete non-formalized proofs of the following reusable facts.

\begin{itemize}
\item Every \emph{fixed} bounded formula over primitive-recursive atomic predicates has a primitive-recursive Boolean characteristic function. Its truth is equivalent in iPRA to the quantifier-free equation that the characteristic function equals one.
\item External parameters, including precision, should be fixed before induction and universally quantified afterward. This avoids putting a quantifier over all precisions inside an induction predicate.
\item A finite contract DAG has a sound reverse demand pass. At fan-out, demands combine by maximum. A single forward bounded trace can check all local contracts.
\item Backward viability on finite candidate clouds is equivalent to the existence of a complete compatible finite path. Choosing the least viable successor is primitive recursive and invokes no infinite dependent choice.
\item If the terminal observable is constant on all viable paths, or has an explicit diameter bound, internal tie-breaking can be erased. This is the right replacement for a discontinuous canonical selector.
\item A partial numerical operation becomes an iPRA construction after totalization and the addition of a positive rational guard. Mere negated equality is not enough for reciprocal, pivoting, chart selection, or branch isolation.
\item Effective compactness is represented by finite nets, a coverage certificate, and a modulus. Approximate minimization is then bounded search.
\item For a decidable coded-point predicate, an explicit height bound yields a complete finite list, and a certified complete finite list yields a height bound by finite maximum. Thus the deep and elementary parts of Diophantine finiteness separate cleanly.
\end{itemize}

\subsection{What remains honestly open}

The framework does not manufacture data that the underlying mathematics does not supply. In particular it does not prove a general effective Faltings or Siegel height bound; it does not turn Selmer classes into rational points on covering spaces; it does not make ineffective constants effective; it does not uniformly choose a sign of an arbitrary real name, a Puiseux branch from arbitrary Cauchy coefficients, or a rotation axis at the identity; and it does not turn metastability into a Cauchy modulus.

These are not implementation omissions of the same kind as a missing substitution lemma. They are producer obligations or computability barriers. The framework's value is that it makes the distinction machine-visible. A packet can now be classified as having
\begin{enumerate}
\item a complete finite verifier,
\item an explicit but unproved producer theorem,
\item a relational output that avoids the producer, or
\item an impossible unrestricted specification that must be reformulated.
\end{enumerate}

\subsection{Deliverables in this archive}

The archive contains:
\begin{itemize}
\item this report and its complete \LaTeX{} source;
\item Lean source blueprints using the exact pinned \code{ConciseIPRA} API;
\item a dependency-ordered Lean completion plan;
\item an exact Python reference kernel implementing bounded folds, trace checking, contract-DAG demand propagation, finite viability, three-way guards, guarded reciprocal, finite-net minimization, height/list conversion, and cubic bisection;
\item eleven passing regression tests and a serialized cubic-inverse certificate;
\item a theorem ledger, a difficulty matrix, a machine-readable framework schema, source provenance, hashes, and PDF validation records.
\end{itemize}

The Python layer is deliberately modest. Its purpose is to expose finite algorithms and catch quantifier-order or certificate-format mistakes. It is not counted as a proof of an iPRA theorem. The Lean files are similarly explicit about their status: they give definitions, exact theorem signatures, and small remaining milestones, but retain visible \code{sorry} placeholders and were not compiled in the build environment.
